% Generated draft: requires mathematical and source review.
\documentclass[11pt]{article}
\usepackage[margin=25mm]{geometry}
\usepackage{fontspec}
\setmainfont{Latin Modern Roman}
\usepackage{amsmath,amssymb,mathtools,unicode-math}
\setmathfont{Latin Modern Math}
\usepackage{enumitem,needspace,xurl,hyperref}
\setlength{\parindent}{0pt}
\setlength{\parskip}{5pt}
\setlength{\emergencystretch}{3em}
\newcommand{\N}{\mathbb{N}}
\newcommand{\Z}{\mathbb{Z}}
\newcommand{\Q}{\mathbb{Q}}
\newcommand{\R}{\mathbb{R}}
\newcommand{\C}{\mathbb{C}}
\newcommand{\F}{\mathbb{F}}
\newcommand{\A}{\mathbb{A}}
\newcommand{\PP}{\mathbb{P}}
\newcommand{\E}{\mathbb{E}}
\newcommand{\eps}{\varepsilon}
\newcommand{\dd}{\,\mathrm{d}}
\begin{document}
\section*{UnsolvedMath ID: 10900027}
\textbf{Source problem code:} catalogue entry
\subsection{Definitions and mathematical statement}
For a closed hyperbolic \(3\)-manifold \(M\), define its global injectivity radius by
\[
 \operatorname{inj}(M)=\inf_{x\in M}\operatorname{inj}_x(M),
\]
equivalently half the length of its shortest closed geodesic. The infimum is over all points; largeness at some point or on most of the volume is a different condition. A fibered \(3\)-manifold here conventionally means one admitting a surface bundle \(M\to S^1\). Homology is interpreted integrally in the two questions, subject to confirming source conventions.

Agol's Problem 5.2 asks two separate existence questions:
\begin{enumerate}
\item For every \(R>0\), does there exist a fibered closed hyperbolic \(3\)-manifold \(M\) such that
\[
 H_i(M;\Z)\cong H_i(S^2\times S^1;\Z)\quad\text{for all }i,
 \qquad\operatorname{inj}(M)>R?
\]
\item For every \(R>0\), does there exist a hyperbolic integral homology \(3\)-sphere \(N\), so \(H_i(N;\Z)\cong H_i(S^3;\Z)\) for all \(i\), with
\[
 \operatorname{inj}(N)>R?
\]
\end{enumerate}
The closed-manifold interpretation is usual for these homology-manifold terms, but is not written explicitly in the first sentence of the source statement and requires confirmation. No fixed fiber genus, common covering family, or uniform fibration is required in the record.

The supplied source mentions related rational-homology-sphere results, and the summary describes them as having bounded rather than arbitrarily large injectivity radius. It records no resolution of either strengthening and reports no verified construction with arbitrarily large global injectivity radius in the stated integral classes. These are supplied status qualifications, not independently checked conclusions.
\subsection{Short English statement}
Can fibered hyperbolic integral homology \(S^2\times S^1\) manifolds have arbitrarily large global injectivity radius? Ask the same for hyperbolic integral homology spheres, without the fibration requirement. Each question requires examples exceeding every positive radius threshold.
\subsection{Sources}
\begin{itemize}
\item Local UnsolvedMath registry, permanent ID \(10900027\), exact source code \texttt{catalogue entry}, subsection 5.2, attributed to Agol; no proposed year supplied.
\item Delp, Hoffoss and Manning, \emph{Problems In Groups, Geometry, and Three-Manifolds} (2015), 5.2, PDF page 4: \url{https://arxiv.org/pdf/1512.04620#page=4}; supplied status-source URL \url{https://arxiv.org/abs/1512.04620}.
\item The source cites related rational-homology-sphere results but no bibliographic details are supplied in the record; it records no resolution of either stated question. The summary reports no verified primary-literature construction for the arbitrarily-large-radius strengthenings.
\item Source-TeX extraction checked against rendered PDF text, accessed 2026-07-29. Source evidence NEEDS\_REVIEW: open in 2015, with no complete modern primary-literature audit recorded. Registry status open and classification OPEN-TRIAGE; summary says likely open, without a verified resolution.
\item Dataset attribution: Ulam AI / UnsolvedMath Contributors, CC BY 4.0. Source rights NEEDS\_REVIEW: publicly accessible author-submitted source TeX and PDF, redistribution terms requiring release review.
\end{itemize}
\end{document}
